\documentclass[a4paper,12pt]{article}
\usepackage[T2A]{fontenc}
\usepackage[utf8]{inputenc}
\usepackage[russian]{babel}
\usepackage{amsmath,amssymb,mathtools}
\usepackage{geometry}
\usepackage{helvet}
\renewcommand{\familydefault}{\sfdefault}
\usepackage{microtype}
\usepackage{booktabs,tabularx,array}
\usepackage{enumitem}
\usepackage{xcolor}
\usepackage{tikz}
\usepackage{hyperref}
\usepackage{fancyhdr}
\usepackage{parskip}
\geometry{left=22mm,right=22mm,top=22mm,bottom=24mm}
\definecolor{accent}{HTML}{007D80}
\definecolor{accentdark}{HTML}{00575C}
\definecolor{ink}{HTML}{242D34}
\definecolor{softgray}{HTML}{69747D}
\definecolor{warn}{HTML}{B02935}
\hypersetup{colorlinks=true,linkcolor=accentdark,urlcolor=accentdark,pdftitle={Ревью домашней работы — 12 ноября 2026 года}}
\setlength{\parindent}{0pt}
\setlength{\parskip}{6pt}
\setlist[itemize]{topsep=2pt,itemsep=3pt,leftmargin=6mm}
\pagestyle{fancy}
\fancyhf{}
\fancyhead[L]{\footnotesize\color{softgray} Математика / проверка домашней работы}
\fancyhead[R]{\footnotesize\color{softgray} 12.11.26}
\fancyfoot[C]{\footnotesize\color{softgray}\thepage}
\renewcommand{\headrulewidth}{0.3pt}
\renewcommand{\footrulewidth}{0pt}
\newcommand{\smallheading}[1]{\par\vspace{5pt}{\color{accentdark}\bfseries #1}\par\vspace{-1pt}}
\newcommand{\markerror}{\textcolor{warn}{\bfseries Ошибка:}}
\newcommand{\markidea}{\textcolor{blue!60!black}{\bfseries Ключевая идея:}}
\newcommand{\markresult}{\textcolor{teal!60!black}{\bfseries Итог:}}
\begin{document}
\thispagestyle{empty}
\begin{center}
\vspace*{37mm}
{\large\color{softgray} МАТЕМАТИКА}\par
\vspace{18mm}
{\LARGE\bfseries\color{ink} Ревью домашней работы}\par
\vspace{7mm}
{\large\color{accentdark} Подробный разбор ошибок}\par
\vspace{31mm}
{\large Лёвин Артём Александрович}\par
\vspace{5mm}
{\large 12 ноября 2026 года}\par
\end{center}
\vfill
\begin{center}
\begin{tikzpicture}[x=1cm,y=1cm]
\draw[accent!65,line width=0.85pt] (-6,0) .. controls (-3,0.65) and (-1,-0.45) .. (1,0.08)
  .. controls (3,0.6) and (5,-0.26) .. (6,0);
\end{tikzpicture}
\end{center}
\clearpage

\section*{Сводная таблица}
В таблицу включены только задания, в решениях которых обнаружены ошибки или незавершённые шаги. Номер задания ведёт к подробному разбору.
\par\vspace{8pt}
\renewcommand{\arraystretch}{1.4}
\small
\par\noindent\begin{tabularx}{\linewidth}{@{}p{11mm}p{24mm}X p{26mm}p{23mm}@{}}
\toprule
\textbf{№} & \textbf{Тема} & \textbf{Что исправить} & \textbf{Ваш ответ} & \textbf{Верный ответ}\\
\midrule
\hyperref[task:2]{\textbf{2}} & Синус & Не завершено вычисление квадрата синуса; результат выражения не получен. & Не указан & $-2$\\
\hyperref[task:4]{\textbf{4}} & Двойной угол & При преобразовании потеряны скобки и неверно записано распределение множителя $30$. & $-\dfrac{690}{25}$ & $-\dfrac{690}{25}$\\
\bottomrule
\end{tabularx}\par
\normalsize
\vfill
{\footnotesize\color{softgray} Важное замечание: совпадение окончательного ответа с верным значением ещё не подтверждает корректность всех промежуточных равенств.}
\clearpage

\thispagestyle{empty}
\vspace*{0.27\textheight}
\begin{center}
{\LARGE\bfseries\color{accentdark} Разбор ошибок}\par
\vspace{8mm}
{\large\color{softgray} Пошаговое восстановление метода}\par
\end{center}
\clearpage

\phantomsection\label{task:2}
\section*{Задание 2. Квадрат синуса и радикалы}
\smallheading{Текст задания}
Найдите значение выражения
\[
\sqrt8-\sqrt{32}\,\sin^2\!\frac{11\pi}{8}.
\]
\smallheading{Ваше решение}
В записи выполнены переходы к степеням двойки:
\[
\sqrt8-\sqrt{32}\,\sin^2\!\frac{11\pi}{8}
=2^{3/2}-2^{5/2}\sin^2\!\frac{11\pi}{8}.
\]
Затем вычисление обрывается. \textbf{Ваш ответ:} не указан.

\smallheading{Где решение остановилось}
Последний корректный шаг: замена квадратных корней дробными степенями.
Первый недостающий шаг: вычисление \(\sin^2(11\pi/8)\). Без этого нельзя завершить численное упрощение.

\smallheading{Почему этого шага не хватает}
\markerror\ В полученном выражении по-прежнему присутствует тригонометрическая функция. Запись \(2^{3/2}-2^{5/2}\sin^2(11\pi/8)\) равносильна исходной, однако ещё не является числовым ответом. Нужно использовать формулу понижения степени и точное значение косинуса.

\markidea\ Для любого угла \(u\)
\[
\sin^2 u=\frac{1-\cos 2u}{2}.
\]
Так квадрат синуса сводится к косинусу удвоенного угла, значение которого здесь известно по единичной окружности.

\smallheading{Исправление от последнего правильного шага}
\[
\begin{aligned}
\sin^2\!\frac{11\pi}{8}
&=\frac{1-\cos(11\pi/4)}{2}\\
&=\frac{1-\cos(3\pi/4)}{2}\\
&=\frac{1-(-\sqrt2/2)}{2}
=\frac{2+\sqrt2}{4}.
\end{aligned}
\]
Здесь использована периодичность: \(11\pi/4=2\pi+3\pi/4\), а \(\cos(3\pi/4)=-\sqrt2/2\).

\smallheading{Полное правильное решение}
Сначала упрощаем радикалы:
\[
\sqrt8=2\sqrt2,\qquad \sqrt{32}=4\sqrt2.
\]
Подставляем значение квадрата синуса:
\[
\begin{aligned}
\sqrt8-\sqrt{32}\sin^2\!\frac{11\pi}{8}
&=2\sqrt2-4\sqrt2\cdot\frac{2+\sqrt2}{4}\\
&=2\sqrt2-\sqrt2(2+\sqrt2)\\
&=2\sqrt2-2\sqrt2-2\\
&=-2.
\end{aligned}
\]
\smallheading{Проверка}
\(\sin^2(11\pi/8)=(2+\sqrt2)/4\) лежит между нулём и единицей, как и должно быть. После подстановки слагаемые с \(\sqrt2\) взаимно уничтожаются.

\markresult\ Значение исходного выражения: \(\boxed{-2}\).

\smallheading{Задача на закрепление}
Вычислите \(\displaystyle\sqrt{18}-\sqrt{72}\,\sin^2\!\frac{7\pi}{8}\).
\clearpage

\phantomsection\label{task:4}
\section*{Задание 4. Формула косинуса двойного угла}
\smallheading{Текст задания}
Найдите \(30\cos 2\alpha\), если \(\cos\alpha=\dfrac15\).

\smallheading{Ваше решение и Ваш ответ}
Верно найдены
\[
\sin^2\alpha=1-\frac1{25}=\frac{24}{25}.
\]
Однако в промежуточной записи множитель \(30\) стоит только перед \(\cos^2\alpha\), а затем появляется переход вида
\[
30\cdot\frac1{25}-\frac{24}{25}
\overset{\text{ошибочно}}{=}
30\left(-\frac{23}{25}\right).
\]
Записанный окончательный ответ: \(\displaystyle-\frac{690}{25}\). Он совпадает с верным значением, но указанное равенство ошибочно.

\smallheading{Где именно возникает ошибка}
Последний правильный шаг: \(\sin^2\alpha=\frac{24}{25}\).
Первый неправильный шаг: потеря скобок после числа \(30\) и некорректное равенство между выражениями с разным количеством множителей \(30\).

\smallheading{Почему это неверно}
\markerror\ По формуле двойного угла
\[
30\cos 2\alpha=30\bigl(\cos^2\alpha-\sin^2\alpha\bigr).
\]
Множитель \(30\) относится \textbf{ко всей разности}. Раскрывая скобки, получаем
\[
30\cos^2\alpha-30\sin^2\alpha.
\]
Если написать \(30\cos^2\alpha-\sin^2\alpha\), второй член уменьшится в тридцать раз; это уже другое выражение. Действительно,
\[
30\cdot\frac1{25}-\frac{24}{25}
=\frac6{25},\qquad
30\left(\frac1{25}-\frac{24}{25}\right)
=-\frac{690}{25}.
\]
Эти числа различны, поэтому переход между ними недопустим.

\markidea\ Сохраняйте скобки, пока не выполнили все действия внутри них. Можно также пользоваться формулой \(\cos 2\alpha=2\cos^2\alpha-1\), которая требует меньше вычислений.

\smallheading{Исправление от последнего правильного шага}
\[
\begin{aligned}
30\cos 2\alpha
&=30\bigl(\cos^2\alpha-\sin^2\alpha\bigr)\\
&=30\left(\frac1{25}-\frac{24}{25}\right)\\
&=30\left(-\frac{23}{25}\right)
=-\frac{690}{25}.
\end{aligned}
\]
\smallheading{Полное правильное решение}
Выберем короткую формулу и последовательно подставим данные:
\[
\begin{aligned}
\cos 2\alpha&=2\cos^2\alpha-1\\
&=2\left(\frac15\right)^2-1\\
&=\frac2{25}-\frac{25}{25}=-\frac{23}{25},\\[2pt]
30\cos 2\alpha&=30\cdot\left(-\frac{23}{25}\right)
=-\frac{690}{25}=-\frac{138}{5}.
\end{aligned}
\]
\smallheading{Проверка}
Так как \(\cos\alpha=\frac15\), косинус двойного угла отрицателен: \(\cos 2\alpha=-\frac{23}{25}\). Это согласуется с найденным знаком. Дополнительно \(\left|-\frac{23}{25}\right|\leq1\), что необходимо для косинуса.

\markresult\ Правильный ответ: \(\boxed{-\dfrac{690}{25}}\). Исправить требуется способ записи промежуточных преобразований.

\smallheading{Задача на закрепление}
Найдите \(18\cos 2\beta\), если \(\cos\beta=\dfrac13\).

\vfill
\begin{center}
{\footnotesize\color{softgray} Лёвин Артём Александрович эксклюзивно для Марка Гукина}
\end{center}
\end{document}
